\documentclass[../../../include/open-logic-section]{subfiles}

\begin{document}

\olfileid[pt]{sth}{spine}{Valphabasic}
\olsection{Propriedades básicas dos estágios}

Para evidenciar a importância fundacional da definição dos
$V_\alpha$, apresentaremos alguns resultados básicos sobre eles. Começamos com uma definição:\footnote{Não há terminologia padrão para «potente»; este é o nome usado por \citet{ButtonLT1}.}
\begin{defn}
	O conjunto $A$ é \emph{potente} sse $\forall x((\exists y \in A)x \subseteq y \lif x \in A)$. 
\end{defn}
\begin{lem}\ollabel{Valphabasicprops}
Para cada ordinal $\alpha$:
\begin{enumerate}
	\item\ollabel{Valphatrans} Cada $V_\alpha$ é transitivo.
	\item\ollabel{Valphapotent} Cada $V_\alpha$ é potente.
	\item\ollabel{Valphacum} Se $\gamma \in \alpha$, então $V_\gamma
	\in V_\alpha$ (e portanto também $V_\gamma \subseteq V_\alpha$, por
	\olref{Valphatrans})
\end{enumerate}
\end{lem}

\begin{proof}
Provamos isto por indução transfinita (simultânea). Para a
indução, supomos que \olref{Valphatrans}--\olref{Valphacum} valem
para cada ordinal $\beta < \alpha$. 

O caso $\alpha = \emptyset$ é trivial. 

Suponhamos $\alpha = \ordsucc{\beta}$. Para mostrar \olref{Valphacum}, se
$\gamma \in \alpha$ então $V_\gamma \subseteq V_\beta$ por hipótese,
logo $V_\gamma \in \Pow{V_\beta} = V_\alpha$. Para mostrar
\olref{Valphapotent}, suponhamos $A \subseteq B \in V_\alpha$, O.s.\@, $A
\subseteq B \subseteq V_\beta$; então $A \subseteq V_\beta$, logo $A \in
V_\alpha$. Para mostrar \olref{Valphatrans}, notemos que se $x \in A \in
V_\alpha$ temos $A \subseteq V_\beta$, logo $x \in V_\beta$, logo $x
\subseteq V_\beta$ pois $V_\beta$ é transitivo por hipótese, e assim $x
\in V_\alpha$. 

Suponhamos que $\alpha$ é um ordinal limite. Para mostrar \olref{Valphacum}, se
$\gamma \in \alpha$ então $\gamma \in \ordsucc{\gamma} \in \alpha$, de modo que
$V_\gamma \in V_{\ordsucc{\gamma}}$ por hipótese, logo
$V_\gamma \in \bigcup_{\beta \in \alpha} V_\beta = V_\alpha$. Para mostrar
\olref{Valphatrans} e \olref{Valphapotent}, basta observar que uma
união de conjuntos transitivos (respetivamente, potentes) é transitiva
(respetivamente, potente). 
\end{proof}

\begin{lem}\ollabel{Valphanotref}
Para cada ordinal $\alpha$, $V_\alpha \notin V_\alpha$.
\end{lem}

\begin{proof}
Por indução transfinita. Evidentemente $V_\emptyset \notin V_\emptyset$. 

Se $V_{\ordsucc{\alpha}} \in V_{\ordsucc{\alpha}} = \Pow{V_\alpha}$,
então $V_{\ordsucc{\alpha}} \subseteq V_\alpha$; e como $V_\alpha
\in V_{\ordsucc{\alpha}}$ por \olref{Valphabasicprops}, temos
$V_\alpha \in V_\alpha$. Reciprocamente: se $V_\alpha \notin V_\alpha$
então $V_{\ordsucc{\alpha}} \notin V_{\ordsucc{\alpha}}$

Se $\alpha$ é limite e $V_\alpha \in V_\alpha = \bigcup_{\beta \in
\alpha}V_\beta$, então $V_\alpha \in V_\beta$ para algum $\beta \in
\alpha$; mas então também $V_\beta \in V_\alpha$, logo $V_\beta \in
V_\beta$ por \olref{Valphabasicprops} (duas vezes). Reciprocamente, se $V_\beta
\notin V_\beta$ para todo o $\beta \in \alpha$, então $V_\alpha \notin
V_\alpha$.
\end{proof}

\begin{cor}
Para quaisquer ordinais $\alpha, \beta$: $\alpha \in \beta$ sse $V_\alpha \in V_\beta$
\end{cor}

\begin{proof}
\Olref{Valphabasicprops} dá um sentido. Reciprocamente, suponhamos $V_\alpha \in V_\beta$. Então $\alpha \neq \beta$ por \olref{Valphanotref}; e $\beta \notin \alpha$, pois caso contrário teríamos $V_\beta \in V_\alpha$ e logo $V_\beta \in V_\beta$ por \olref{Valphabasicprops} (duas vezes), contradizendo \olref{Valphanotref}. Logo $\alpha \in \beta$ por tricotomia.
\end{proof}
\noindent
Tudo isto permite pensar em cada $V_\alpha$ como o $\alpha$-ésimo
estágio da hierarquia. Eis a razão.

Certamente os nossos $V_\alpha$ podem ser vistos como formados num
processo \emph{iterativo}, pois o uso de ordinais acompanha a noção de
iteração. Além disso, se um estágio é formado antes do outro, O.s.\@,
$V_\beta \in V_\alpha$, O.s.\@, $\beta \in \alpha$, então o nosso processo de
formação é \emph{cumulativo}, pois $V_\beta \subseteq V_\alpha$.
Finalmente, estamos de fato a formar \emph{todas} as coleções possíveis de conjuntos
que estavam disponíveis em qualquer estágio anterior, pois qualquer estágio sucessor
$V_{\ordsucc{\alpha}}$ é o conjunto potência do seu predecessor $V_\alpha$.

Em resumo: com $\ZFminus$, estamos \emph{quase} a terminar de articular
a nossa visão da hierarquia cumulativa e iterativa de conjuntos. (Embora,
naturalmente, ainda precisemos de justificar o Esquema de Substituição.)

\end{document}
